\clearpage\thispagestyle{otFront}\phantomsection\label{nouveautes}\pdfbookmark[0]{Nouveau programme}{bm-nouveautes}\otPageTitre{Nouveau programme}{Ce qui change en Troisième à la rentrée 2028}\fontsize{11}{15}\selectfont

Le programme de mathématiques du cycle 4 publié en 2026 s’applique en Troisième
à la rentrée 2028. Cette édition a été réécrite en conséquence ; voici les
principaux changements par rapport au programme précédent.\par\medskip
\fontsize{10}{13.5}\selectfont
\begin{tabularx}{\linewidth}{@{}>{\raggedright\arraybackslash}p{27mm}Y@{}}\toprule
Nouveau en Troisième & Translations et vecteurs : vecteurs égaux, somme, relation de Chasles (ch. 15). Effectifs cumulés, quartiles et boîtes à moustaches (ch. 10). Réunion, intersection et relation $P(A\cup B)+P(A\cap B)=P(A)+P(B)$ (ch. 11). Inéquations $ax\ge b$ par le calcul et graphiquement (ch. 5 et 9). Trois identités remarquables et raisonnement par analyse-synthèse (ch. 4 et 5). Fonction carré et équation $x^2=a$ (ch. 2 et 8). Sphère, boule, grands cercles et sections (ch. 16). Conditions composées, boucles conditionnelles et blocs personnalisés (ch. 17).\\[3mm]
Prolongements & Notations $\mathbb N$, $\mathbb Z$, $\mathbb D$, $\mathbb Q$ (ch. 1) ; préfixes de nano à giga (ch. 2) ; identité de Sophie Germain et cylindre d’Archimède (défis et ch. 16).\\[3mm]
Réactivé & Pythagore, sa réciproque et sa contraposée, cercle circonscrit au triangle rectangle, droite des milieux, grandeurs quotients et expériences à deux épreuves : étudiés en Quatrième, ils sont attendus en automatismes et dans les problèmes (ch. 7, 11 et 12).\\[3mm]
Retiré & Homothéties et rotations, triangles semblables, effet d’un agrandissement sur les aires et les volumes, repérage sur la sphère et dans un pavé, histogrammes : ces notions ne figurent plus dans le programme du cycle 4 et ne sont plus traitées.\\\bottomrule
\end{tabularx}
\par\medskip\fontsize{11}{15}\selectfont
Le vocabulaire suit le programme : quotient, fraction et nombre rationnel ;
demi-tour pour la symétrie centrale ; théorème, réciproque et contraposée ;
évènement, réunion et intersection. La notation $x\mapsto f(x)$ est utilisée
pour les fonctions. La matrice de fin de livre relie chaque rubrique du
programme aux chapitres et aux exercices.

\clearpage
